\subsection{Functions}

A function is a constraint imposed on a mathematical relationship between
a set of inputs and a set of outputs. Only two properties
are required to be satisfied for a relatioship to be called a function. Each
input must have an output and each output must be unique.

Each input can be thought of as a dimension in an n-dimensional space where
n is the total number of inputs and thus a function forms a n-dimensional
surface in that space representing all possible outputs.

Here are some example functions:
\begin{align*}
  \mathbf{F}(x, y) = [f_1(x, y), f_2(x, y), \ldots, f_n(x, y)] = \text{vector output} \\
  \mathbf{F}(x, y, z) = [f_1(x, y, z), f_2(x, y, z), \ldots, f_n(x, y, z)] = \text{vector output} \\
  \mathbf{F}(x) =  12x^2 + 3x + 5 = \text{scalar output} \\
  \mathbf{G}(x, y) = x^2 + y^2 = \text{scalar output} \\
  \mathbf{H}(x, y, z) = \begin{bmatrix}
    x + y + z & 2x - y + z & x - 2y + 3z \\
    x - y + z & 3x + y - z & 2x + 2y + z \\
    x + y - z & x - 3y + 2z & 4x + y + z
  \end{bmatrix} = \text{tensor output}
\end{align*}

\subsubsection{Classification}

Functions are classified based on the number of output values they produce.
\begin{enumerate}
  \item Scalar Function: The output is a single value that can be
    represented by a single number. Temperature, pressure, and density
    are examples of scalar functions. Another example,
    \begin{align*}
      F(x, y, z) = 20 + 0.5x - 0.2y + 0.1z
    \end{align*}
  \item Vector Function: The output is a vector that can be represented
    by multiple numbers. Wind velocity, electric function, and magnetic
    function are examples of vector functions. Another example,
    \begin{align*}
      \mathbf{F}(x, y, z) = \begin{bmatrix}
      2x - y + z \\
      x + 3y - 2z \\
      -x + y + 4z
      \end{bmatrix}
    \end{align*}
  \item Tensor Function: The output is a tensor that can be represented
    by a multidimensional array of numbers. Stress, strain, and moment
    of inertia are examples of tensor functions. Another example,
    \begin{align*}
      \mathbf{F}(x, y, z) = \begin{bmatrix}
      2x - y + z & x + 3y - 2z & -x + y + 4z \\
      3x - 2y + z & 2x + y - z & -x + 2y + 3z \\
      -2x + y + z & x - 3y + 2z & 4x + y - z
      \end{bmatrix}
    \end{align*}
\end{enumerate}

\subsubsection{Notation}

Scalar functions are represented by normal font and Vector functions are
represented by boldface font.
\begin{align*}
  V(x, y, z) &\implies \text{Scalar Function} \
  \mathbf{F}(x, y, z) &\implies \text{Vector Function}
\end{align*}
\
Often times \( \mathbf{r} \) is used to represent the spatial coordinates
such that,
\begin{align*}
  \mathbf{r} &= (x, y, z) \\
  F(\mathbf{r}) &= F(x) \hspace{1cm} \text{[where r=(x); 1D function]} \\
  F(\mathbf{r}, t) &= F(x, y, z, t) \hspace{1cm} \text{[where r=(x, y, z); 3D function]} \\
  G(\mathbf{r}) &= G(a, b) \hspace{1cm} \text{[where r=(a, b); 2D function]}
\end{align*}
\
We \textbf{will refrain from abstracting/hiding information for ease of
understanding}.

\subsubsection{Operations}

In functions of physics, the inputs are mostly spatial coordinates and
time i.e. \((x, y, z, t)\) while the outputs are mostly scalars or
vectors. We will thus work with scalar and vector functions. For a scalar
function, operations are performed pointwise (f at a, b can only go with
g at a, b). For a vector function, operations are performed component-wise
(f at x goes with g at x, f at y goes with g at y, f at z goes with g at
z{\ldots}so on).
\
Assume the following functions exist:
\begin{align*}
  V(x, y, z) \hspace{1cm} \text{(i) Scalar Function} \\
  W(x, y, z) \hspace{1cm} \text{(ii) Scalar Function} \\
  \mathbf{F}(x, y, z) = \left[F_x(x, y, z), F_y(x, y, z), F_z(x, y, z)\right] \hspace{1cm} \text{(iii) Vector Function} \\
  \mathbf{G}(x, y, z) = \left[G_x(x, y, z), G_y(x, y, z), G_z(x, y, z)\right] \hspace{1cm} \text{(iv) Vector Function}
\end{align*}
\textbf{How many ways do you think we can combine V, W, F and G}? There
are three possibilities if we assume commutativity:
\begin{enumerate}
  \item Scalar on Scalar: \(V(x, y, z)\) \& \(W(x, y, z)\) |or| \(W(x, y, z)\) \& \(V(x, y, z)\)
  \item Scalar on Vector: \(V(x, y, z)\) \& \(\mathbf{F}(x, y, z)\) |or| \(\mathbf{F}(x, y, z)\) \& \(V(x, y, z)\)
  \item Vector on Vector: \(\mathbf{F}(x, y, z)\) \& \(\mathbf{G}(x, y, z)\) |or| \(\mathbf{G}(x, y, z)\) \& \(\mathbf{F}(x, y, z)\)
\end{enumerate}
\
Let's try arithmetic operations on these combos.
\begin{enumerate}
  \item \textbf{Addition and Subtraction:} This is straightforward. The
    result is a new function of the same type as the input functions.
    \begin{enumerate}
      \item Scalar on Scalar:
        \begin{align*}
          R(x, y, z) &= V(x, y, z) + W(x, y, z) \\
          R(x, y, z) &= V(x, y, z) - W(x, y, z)
        \end{align*}
      \item Scalar on Vector:
        \begin{align*}
          \mathbf{R}(x, y, z) &= V(x, y, z) + \mathbf{F}(x, y, z) \\
                               &= \begin{bmatrix}
                                    V(x, y, z) + F_x(x, y, z) \\
                                    V(x, y, z) + F_y(x, y, z) \\
                                    V(x, y, z) + F_z(x, y, z)
                                  \end{bmatrix} \\
          \mathbf{R}(x, y, z) &= V(x, y, z) - \mathbf{F}(x, y, z) \\
                               &= \begin{bmatrix}
                                    V(x, y, z) - F_x(x, y, z) \\
                                    V(x, y, z) - F_y(x, y, z) \\
                                    V(x, y, z) - F_z(x, y, z)
                                  \end{bmatrix}
        \end{align*}
      \item Vector on Vector:
        \begin{align*}
          \mathbf{R}(x, y, z) &= \mathbf{F}(x, y, z) + \mathbf{G}(x, y, z) \\
                               &= \begin{bmatrix}
                                    F_x(x, y, z) + G_x(x, y, z) \\
                                    F_y(x, y, z) + G_y(x, y, z) \\
                                    F_z(x, y, z) + G_z(x, y, z)
                                  \end{bmatrix} \\
          \mathbf{R}(x, y, z) &= \mathbf{F}(x, y, z) - \mathbf{G}(x, y, z) \\
                               &= \begin{bmatrix}
                                    F_x(x, y, z) - G_x(x, y, z) \\
                                    F_y(x, y, z) - G_y(x, y, z) \\
                                    F_z(x, y, z) - G_z(x, y, z)
                                  \end{bmatrix}
        \end{align*}
    \end{enumerate}
  \item \textbf{Multiplication:} More work is required as the kind of
    functions being multiplied now matters. The result is a function
    of the same or different type.
    \begin{enumerate}
    \item Scalar on Scalar: Results in a scalar function.
      \begin{align*}
      R(x, y, z) = V(x, y, z) \cdot W(x, y, z)
      \end{align*}
    \item Scalar on Vector: Results in a vector function.
      \begin{align*}
        \mathbf{R}(x, y, z) &= V(x, y, z) \cdot \mathbf{F}(x, y, z) \\
                             &= \begin{bmatrix}
                                  V(x, y, z) \cdot F_x(x, y, z) \\
                                  V(x, y, z) \cdot F_y(x, y, z) \\
                                  V(x, y, z) \cdot F_z(x, y, z)
                                \end{bmatrix} \\
        R(x, y, z) &= \mathbf{F}(x, y, z) \cdot V(x, y, z) \\
                             &= \begin{bmatrix}
                                  F_x(x, y, z) \cdot V(x, y, z) \\
                                  F_y(x, y, z) \cdot V(x, y, z) \\
                                  F_z(x, y, z) \cdot V(x, y, z)
                                \end{bmatrix}
      \end{align*}
    \item Vector on Vector: Results in either a scalar function
      (dot product) or a vector function (cross product). \textbf{This
      is worth reading about in more detail. It is used extensively
      in electromagnetism.}
      \begin{align*}
        R(x, y, z) &= \mathbf{F}(x, y, z) \cdot \mathbf{G}(x, y, z) \\
                   &= F_x G_x + F_y G_y + F_z G_z
                   \hspace{1em} \text{[Dot Product]} \\
        \mathbf{R}(x, y, z) &= \mathbf{F}(x, y, z) \times \mathbf{G}(x, y, z) \\
                             &= \begin{bmatrix}
                                  F_y G_z - F_z G_y \\
                                  F_z G_x - F_x G_z \\
                                  F_x G_y - F_y G_x
                                \end{bmatrix}
                   \hspace{1em} \text{[Cross Product]}
      \end{align*}
    \end{enumerate}
\end{enumerate}

\subsubsection{Del/Nabla}
Del (\(\nabla\)) is a vector function. It is defined as:
\begin{align*}
  \nabla = \left[\frac{\partial}{\partial x}, \frac{\partial}{\partial y}, \frac{\partial}{\partial z}\right] = \hat{\mathbf{x}} \frac{\partial}{\partial x} + \hat{\mathbf{y}} \frac{\partial}{\partial y} + \hat{\mathbf{z}} \frac{\partial}{\partial z} \hspace{1em} \text{[3D]} \\
  \nabla = \left[\frac{\partial}{\partial x}, \frac{\partial}{\partial y}\right] = \hat{\mathbf{x}} \frac{\partial}{\partial x} + \hat{\mathbf{y}} \frac{\partial}{\partial y} \hspace{1em} \text{[2D]}
\end{align*}

Think of it as just another vector function that can be combined with other
functions. \textbf{At heart, nabla obeys the same rules of vector
multiplication}.
\begin{itemize}
  \item \textbf{Gradient}: Acts on a scalar function \( V(x, y, z) \) to produce
    a vector function. Intuitively, the \textbf{gradient points in the direction of
    the greatest rate of increase of the scalar function and its magnitude
    represents the rate of change in that direction}. It is defined as:
    \begin{align*}
      \nabla V(x,y,z) = \left[\frac{\partial V(x,y,z)}{\partial x}, \frac{\partial V(x,y,z)}{\partial y}, \frac{\partial V(x,y,z)}{\partial z}\right]
    \end{align*}
  \item \textbf{Divergence}:  Acts on a vector function \( \mathbf{F}(x, y, z) \)
    to produce a scalar function. Intuitively, \textbf{the divergence
    measures  the \textbf{outflow} (positive) or \textbf{inflow} (negative) of the
    vector function at a given point}. It is defined as:
    \begin{align*}
      \nabla \cdot \mathbf{F}(x,y,z) = \frac{\partial F(x,y,z)}{\partial x} +
      \frac{\partial F(x,y,z)}{\partial y} +
      \frac{\partial F(x,y,z)}{\partial z}
    \end{align*}
  \item \textbf{Curl}: Acts on vector function \( \mathbf{F}(x, y, z) \) to
    produce a new vector function that is orthogonal/normal/perpendicular to
    the original vector function at each point in space. Intuitively, \textbf{
    the curl measures the rotation or circulation of the vector function at a
    given point}. It is defined as:
    \begin{align*}
      \nabla \times \mathbf{F}(x,y,z) &=
      \begin{bmatrix}
        \hat{\mathbf{x}} & \hat{\mathbf{y}} & \hat{\mathbf{z}}\\
        \frac{\partial}{\partial x} &
        \frac{\partial}{\partial y} &
        \frac{\partial}{\partial z}\\
        F(x,y,z) & F(x,y,z) & F(x,y,z)
      \end{bmatrix}\\
      &= \left[\frac{\partial F_z(x,y,z)}{\partial y} - \frac{\partial F_y(x,y,z)}{\partial z}, \frac{\partial F_x(x,y,z)}{\partial z} - \frac{\partial F_z(x,y,z)}{\partial x}, \frac{\partial F_y(x,y,z)}{\partial x} - \frac{\partial F_x(x,y,z)}{\partial y}\right]
    \end{align*}
\end{itemize}

\subsubsection{Borrowing Vector Calculus Results}
Several results from vector calculus are useful. Here are
a few that we will use further down the line.
\begin{enumerate}
  \item Curl of Gradient \(\nabla \times (\nabla V)\) is always zero. Here's the proof:
    \begin{align*}
      \nabla \times (\nabla V) &= \begin{bmatrix}
        \hat{\mathbf{x}} & \hat{\mathbf{y}} & \hat{\mathbf{z}}\\
        \frac{\partial}{\partial x} &
        \frac{\partial}{\partial y} &
        \frac{\partial}{\partial z}\\
        \frac{\partial V}{\partial x} & \frac{\partial V}{\partial y} & \frac{\partial V}{\partial z}
      \end{bmatrix}\\
      &= \left[\frac{\partial^2 V}{\partial y \partial z} - \frac{\partial^2 V}{\partial z \partial y}, \frac{\partial^2 V}{\partial z \partial x} - \frac{\partial^2 V}{\partial x \partial z}, \frac{\partial^2 V}{\partial x \partial y} - \frac{\partial^2 V}{\partial y \partial x}\right] = 0
    \end{align*}
  \item Divergence of Curl \(\nabla \cdot (\nabla \times \mathbf{F})\) is always zero. Here's
  the proof:
      \begin{align*}
        \nabla \cdot (\nabla \times \mathbf{F}) &= \frac{\partial}{\partial x} \left(\frac{\partial F_z}{\partial y} - \frac{\partial F_y}{\partial z}\right) + \frac{\partial}{\partial y} \left(\frac{\partial F_x}{\partial z} - \frac{\partial F_z}{\partial x}\right) + \frac{\partial}{\partial z} \left(\frac{\partial F_y}{\partial x} - \frac{\partial F_x}{\partial y}\right)\\
        &= \frac{\partial^2 F_z}{\partial x \partial y} - \frac{\partial^2 F_y}{\partial x \partial z} + \frac{\partial^2 F_x}{\partial y \partial z} - \frac{\partial^2 F_z}{\partial y \partial x} + \frac{\partial^2 F_y}{\partial z \partial x} - \frac{\partial^2 F_x}{\partial z \partial y}\\
        &= 0
      \end{align*}
  \item Curl of Curl \(\nabla \times (\nabla \times \mathbf{F})\) can be rewritten as
    \( \nabla(\nabla \cdot \mathbf{F}) - \nabla^2 \mathbf{F}\). Here's the proof:
    \begin{align*}
      \nabla \times (\nabla \times \mathbf{F}) = \begin{bmatrix}
        \hat{\mathbf{x}} & \hat{\mathbf{y}} & \hat{\mathbf{z}}\\
        \frac{\partial}{\partial x} &
        \frac{\partial}{\partial y} &
        \frac{\partial}{\partial z}\\
        \frac{\partial F_z}{\partial y} - \frac{\partial F_y}{\partial z} &
        \frac{\partial F_x}{\partial z} - \frac{\partial F_z}{\partial x} &
        \frac{\partial F_y}{\partial x} - \frac{\partial F_x}{\partial y}
      \end{bmatrix}\\
      = \left[\frac{\partial^2 F_y}{\partial x \partial y} - \frac{\partial^2 F_x}{\partial y^2} + \frac{\partial^2 F_x}{\partial z^2} - \frac{\partial^2 F_z}{\partial x \partial z}, \frac{\partial^2 F_z}{\partial x^2} - \frac{\partial^2 F_x}{\partial x \partial z} + \frac{\partial^2 F_x}{\partial y^2} - \frac{\partial^2 F_y}{\partial y^2}, \frac{\partial^2 F_x}{\partial y^2} - \frac{\partial^2 F_y}{\partial x \partial y} + \frac{\partial^2 F_y}{\partial z^2} - \frac{\partial^2 F_z}{\partial y \partial z}\right]\\
      = [(\nabla(\nabla \cdot \mathbf{F}))_x - (\nabla^2 \mathbf{F})_x, (\nabla(\nabla \cdot \mathbf{F}))_y - (\nabla^2 \mathbf{F})_y, (\nabla(\nabla \cdot \mathbf{F}))_z - (\nabla^2 \mathbf{F})_z]\\
      =  (\nabla(\nabla \cdot \mathbf{F})) - (\nabla^2) (\mathbf{F})
    \end{align*}
\end{enumerate}
